% Experiment 3 — evidence accumulation on the held-out worlds (the Exp-2 sparse-evidence view).
% GENERATED: paper/figures/exp3_per_week.pdf from the OAT eval dumps by
% exp3_training_transfer/dag_family/plot_per_week.py (references recomputed at every prefix on the
% identical episodes). Regenerate with:  python3 plot_per_week.py
\begin{figure*}[t]
\centering
\includegraphics[width=\linewidth]{figures/exp3_per_week.pdf}
\caption{\textbf{Evidence accumulation on the held-out worlds.} The Experiment-3 analogue of
Experiment~2's sparse-evidence curve: each axis is read as a function of the number of weeks of history
$k$ the forecaster has seen, for the untrained base (grey), the family-trained model (teal), and the
structureless control (amber), against references recomputed at every $k$ on the identical episodes.
Both arms are read at the same training step and over the same seeds, so any gap is the arm and not the
sample. We pool only the held-out worlds whose probe response at $h^\star$ \emph{is} the gain itself
(\texttt{direct}, \texttt{two\_news\_feedback}); \texttt{mediators} and \texttt{chain4} are excluded
because their response is an affine image of $g$ (so $\varepsilon$ lives on a different scale) and
because they are unidentifiable at small $k$. Including them would change \emph{what is being
averaged} as $k$ grows, and a flat curve could be mistaken for a flat model. Informativeness $\rho$ rises steadily with evidence, showing that estimates
sharpen as polls arrive. The gap between teal and amber identifies the training effect. If the
structureless control tracks the family arm, the improvement over the untrained base reflects format
and level tracking rather than recovered structure. Curves use
all three paired seeds at the completed endpoint.}
\label{fig:exp3-per-week}
\end{figure*}
